% race_proof.tex -- round 460, serious proof attempt at the selected race.
% OUTCOME PARTIAL: exact spectral reduction + finite certificate;
% source-derived cofinal spectral-overlap estimate remains open.
% Builds with pdflatex.
% Companion machine check: verify_race_proof.py.
% Discovery: experiments/tfpt-discovery/race_proof_probe.py.
% NO RH CLAIM.  NO anti-RH claim.
\documentclass[11pt]{article}

\usepackage[a4paper,margin=2.7cm]{geometry}
\usepackage{amsmath,amssymb,amsthm}
\usepackage{booktabs}
\usepackage{microtype}
\usepackage{xcolor}
\usepackage[colorlinks=true,linkcolor=blue!50!black,urlcolor=blue!50!black]{hyperref}

\newtheorem{theorem}{Theorem}
\newtheorem{lemma}{Lemma}
\newtheorem{proposition}{Proposition}
\theoremstyle{remark}
\newtheorem{remark}{Remark}

\newcommand{\R}{\mathbb{R}}
\newcommand{\ip}[2]{\langle #1,#2\rangle}
\newcommand{\qfull}{q^{\dagger}_{\mathrm{full}}}
\newcommand{\qstar}{q_*}

\title{\bfseries The selected-window race\\[2mm]
\large a full-proof attempt and the surviving spectral-overlap gap}
\author{Stefan Hamann}
\date{August 30, 2026}

\begin{document}
\maketitle

\begin{abstract}\noindent
We attempt to prove that
$\Delta q_{\rm arith}(k)<1-\qstar(k)$ for infinitely many members of
the named Lean-selected sequence
$a_k=2^k$, $r_k=\lfloor\sqrt{k}\rfloor$,
$m_k=k2^{r_k}-1$.
The outcome is \textbf{PARTIAL}.
We prove an exact spectral representation of the race and a genuine
two-band sufficient criterion.  The criterion gives a finite
certificate at $k=5$, but no uniform estimate: the raw resolvent norm
bound grows from $4.82$ to $4991.65$ on $k=5,\ldots,12$, and even the
optimal two-band bound exceeds $1$ from $k=6$ onward.  The fixed-mesh
plateaux cannot supply an infinite subsequence because the
$r$th plateau contains exactly $2r+1$ indices.
The precise unresolved statement is a source-derived estimate for the
spectral measure of the explicit border vector near the saturating
edge $\lambda=1$.  No Riemann-hypothesis claim, and no anti-RH claim,
is made.
\end{abstract}

\medskip\noindent
\textbf{Binding outcome.}\enspace
\textbf{PARTIAL}: spectral reduction proved; norm route refuted;
plateau route finite; spectral-overlap estimate open.
The candidate race lemma is not proved, refuted, or promoted.
No $L^*$ claim.  No $R^\dagger$ claim.  No RH claim.

\section{The exact object}

Fix one selected finite window and its cap $N$.
Use the positive part $\mu$ of the folded signed measure as the
orthogonality measure.  In its orthonormal polynomial basis let
\[
 B_{yi}=\sqrt{\nu_y}\,p_i(y),\qquad
 C=B^{T}B,\qquad
 u=\frac{b}{\sqrt{B_w}},
\]
where $\nu$ is the negative part, $b_i$ is the smooth-border
functional on $p_i$, and $B_w>0$ is the standing budget.
These are exactly the matrices used by
\texttt{deep\_builder\_probe.solve\_qdag}.  The finite augmented
matrix in this frame is
\begin{equation}\label{eq:A}
 \mathcal A=
 \begin{pmatrix}
 I-C&u\\ u^T&1
 \end{pmatrix}.
\end{equation}
The r453/r457 identity gives
\begin{equation}\label{eq:race}
 \qfull=\qstar+\Delta q_{\rm arith},\qquad
 \qfull<1
 \ \Longleftrightarrow\
 \Delta q_{\rm arith}<1-\qstar .
\end{equation}
Equation~\eqref{eq:race} is scalar algebra; it does not by itself
prove either inequality.

\section{What can be proved}

\begin{theorem}[Spectral race identity]\label{thm:spectral}
Assume $C\prec I$.  Let
$Ce_j=\lambda_j e_j$ be an orthonormal eigendecomposition.
Then
\begin{equation}\label{eq:spectral}
 \qfull
 =u^T(I-C)^{-1}u
 =\sum_{j=1}^{N}
   \frac{|\ip{u}{e_j}|^2}{1-\lambda_j}.
\end{equation}
Moreover,
\[
 \mathcal A\succ0
 \quad\Longleftrightarrow\quad
 \qfull<1 .
\]
\end{theorem}

\begin{proof}
Because $C$ is real symmetric, the spectral theorem gives
$C=U\operatorname{diag}(\lambda_j)U^T$ with $U$ orthogonal.
The hypothesis $C\prec I$ implies $1-\lambda_j>0$, hence
\[
 (I-C)^{-1}
 =U\operatorname{diag}\bigl((1-\lambda_j)^{-1}\bigr)U^T.
\]
Taking the quadratic form with $u$ proves~\eqref{eq:spectral}.
For the second assertion take the Schur complement of the positive
leading block $I-C$ in~\eqref{eq:A}:
\[
 \mathcal A/(I-C)
 =1-u^T(I-C)^{-1}u=1-\qfull .
\]
The Schur-complement criterion proves the equivalence.
\end{proof}

\begin{theorem}[Two-band sufficient criterion]\label{thm:twoband}
Under the assumptions of Theorem~\ref{thm:spectral}, fix
$\theta<1$ and let
$P_\theta=\mathbf 1_{(\theta,1)}(C)$.
Writing $\lambda_{\max}=\lambda_{\max}(C)$, one has
\begin{equation}\label{eq:twoband}
\qfull\le
\frac{\|(I-P_\theta)u\|^2}{1-\theta}
+\frac{\|P_\theta u\|^2}{1-\lambda_{\max}}.
\end{equation}
Consequently, strict inequality of the right-hand side with $1$
is a sufficient certificate for the race inequality.
\end{theorem}

\begin{proof}
Split the sum~\eqref{eq:spectral} into eigenvalues
$\lambda_j\le\theta$ and $\lambda_j>\theta$.
On the first set,
$(1-\lambda_j)^{-1}\le(1-\theta)^{-1}$.
On the second,
$(1-\lambda_j)^{-1}\le(1-\lambda_{\max})^{-1}$.
Parseval's identity turns the two sums of squared coefficients into
$\|(I-P_\theta)u\|^2$ and $\|P_\theta u\|^2$.
This proves~\eqref{eq:twoband}; combine it with
Theorem~\ref{thm:spectral} and~\eqref{eq:race}.
\end{proof}

\begin{proposition}[The fixed-mesh blocks are finite]\label{prop:block}
For every integer $r\ge0$,
\[
\{k\in\mathbb N:\lfloor\sqrt{k}\rfloor=r\}
=\{r^2,r^2+1,\ldots,(r+1)^2-1\}
\]
and this set has cardinality $2r+1$.
On it the selected spacing is
\[
\Delta_k=\frac{\log a_k}{m_k+1}
=2^{-r}\log2.
\]
Thus a monotonicity or averaging argument confined to one fixed
$r$ cannot by itself establish an infinitely-often statement.
\end{proposition}

\begin{proof}
The equality $\lfloor\sqrt{k}\rfloor=r$ is equivalent to
$r\le\sqrt{k}<r+1$, hence to
$r^2\le k<(r+1)^2$.  There are
$(r+1)^2-r^2=2r+1$ such integers.
Since $a_k=2^k$ and $m_k+1=k2^r$ on this block,
\[
\Delta_k=\frac{k\log2}{k2^r}=2^{-r}\log2.
\]
Every one of these blocks is finite, which proves the quantifier
statement.
\end{proof}

\section{Numerical adjudication of the three attacks}

The companion probe uses the complete prime-power comb through
$a_k^2=4^k$ and the same fold, border, chain, and solve as r459.
The spectral identity agrees with the direct solve to
$3.85\cdot10^{-14}$ and with the r459 race object to
$1.89\cdot10^{-14}$.

\begin{center}
\small
\begin{tabular}{@{}rrrrrrrr@{}}
\toprule
$k$ & $N$ & race & $\lambda_{\max}(C)$ &
raw bound & best two-band & Gersh. & top overlap\\
\midrule
5  & 10 & .675060 & .8269525 & 4.825 & .910 & $-1.284$ & $9.53\,10^{-3}$\\
6  & 12 & .694907 & .9446612 & 14.227 & 1.319 & $-1.404$ & $2.22\,10^{-6}$\\
7  & 14 & .628736 & .9867839 & 47.207 & 1.167 & $-1.260$ & $8.89\,10^{-4}$\\
8  & 16 & .606566 & .9941763 & 93.728 & 1.726 & $-1.162$ & $5.05\,10^{-4}$\\
9  & 36 & .635369 & .9991328 & 724.586 & 4.382 & $-2.274$ & $1.94\,10^{-5}$\\
10 & 40 & .757783 & .9995864 & 1593.258 & 5.948 & $-2.634$ & $4.15\,10^{-6}$\\
11 & 44 & .715782 & .9998267 & 2755.648 & 6.651 & $-2.122$ & $5.10\,10^{-6}$\\
12 & 48 & .693564 & .9999128 & 4991.649 & 10.831 & $-1.975$ & $1.54\,10^{-6}$\\
\bottomrule
\end{tabular}
\end{center}

\paragraph{Attack 1: norm of the perturbation.}
The canonical one-number resolvent estimate is
\begin{equation}\label{eq:raw}
\qfull\le\frac{\|u\|^2}{1-\lambda_{\max}(C)}.
\end{equation}
It already gives $4.825>1$ at $k=5$ and worsens to
$4991.649$ at $k=12$.
The reason is visible rather than conjectural:
$\lambda_{\max}(C)\uparrow1$, while the squared overlap of $u$ with
the top eigenvector is only
$1.5\cdot10^{-6}$--$9.5\cdot10^{-3}$; the top mode contributes at
most $5.6\%$ of the true $q$ on this census.
The operator norm discards precisely this alignment.
Even optimizing Theorem~\ref{thm:twoband} over every spectral cut
certifies only $k=5$ ($0.910358<1$); its best value is $>1$ on
$k=6,\ldots,12$.

A second obstruction prevents the proposed Weyl comparison with the
ARCH baseline: $\qstar$ is computed at the smaller degree $J_P$,
whereas $\qfull$ is at degree $N$.  Weyl requires two operators on
the same space.  The full-depth ARCH signed chain is not a positive
baseline; it has respectively $3,3,4,4,6,6,7,8$ sign flips on
$k=5,\ldots,12$.  Thus the missing extension from $J_P$ to $N$
would itself contain the positivity problem one hoped to prove.

\paragraph{Attack 2: fixed $r$.}
The ARCH prefix values are numerically identical within the two
observed blocks:
\[
\qstar=0.624231952161870\quad(k=5,\ldots,8),\qquad
\qstar=0.719102723128236\quad(k=9,\ldots,12).
\]
This useful simplification does not address the quantifier:
Proposition~\ref{prop:block} proves that every such block is finite.
An infinite-subsequence proof must compare infinitely many different
$r$ and hence infinitely many different spacings.

\paragraph{Attack 3: direct positivity.}
For every tested main window the true augmented minimum eigenvalue is
positive and $\rho_{N-1}<5/7$, consistently with r459.
But the Gershgorin lower bound of~\eqref{eq:A} lies between
$-2.634$ and $-1.162$ on all eight rows.
Interlacing alone gives no sign because the parent ARCH cap is
indefinite.  Direct positivity therefore returns to the spectral
alignment in Theorem~\ref{thm:spectral}; no source-pure
Markov/interlacing inequality was obtained.

\section{Exact remaining gap}

\noindent\fbox{\parbox{0.94\linewidth}{
\textbf{Remaining gap (verbatim).}
Prove, from the explicit von-Mangoldt tents and smooth-border source
and without evaluating $q_k$ itself, that for infinitely many
selected $k$ one has $C_k\prec I$ and there exists
$\theta_k<1$ such that
\[
\sum_{\lambda_{k,j}>\theta_k}
\frac{|\langle u_k,e_{k,j}\rangle|^2}{1-\lambda_{k,j}}
<
1-
\sum_{\lambda_{k,j}\le\theta_k}
\frac{|\langle u_k,e_{k,j}\rangle|^2}{1-\lambda_{k,j}}.
\]
Equivalently: bound the spectral measure of the explicit normalized
border vector near the saturating edge $\lambda=1$ strongly enough
that its Stieltjes transform at $1$ is $<1$ infinitely often.
No such source-derived overlap estimate is proved here.
}}

\medskip
This statement is deliberately more precise than ``control the
primes'': it identifies the required currency and the quantifier.
It is also honest about its relation to the target: by
Theorem~\ref{thm:spectral} it is exactly the missing inequality,
split into safe and dangerous spectral bands.  The split is useful
only if its two terms can be bounded before the target is read back.

\section{Alias and refutation audit}

\begin{center}
\small
\begin{tabular}{@{}p{3.4cm}p{3.0cm}p{7.0cm}@{}}
\toprule
Route & classification & reason\\
\midrule
Raw norm~\eqref{eq:raw}, Gershgorin &
acknowledged alias / refuted &
Same unsigned-majorant failure class as r387/r429; the probe records
its failure and does not advertise it as a mechanism.\\
Rosser--Schoenfeld or other first-kind $\psi$ upper bounds &
not used to claim progress &
Inserted only through unsigned mass or operator norm, such bounds
erase the overlaps in~\eqref{eq:spectral}; they cannot repair the
measured failure of~\eqref{eq:raw} without additional signed
structure.\\
Fixed-$r$ plateau &
new quantifier audit, not octave additivity &
No increments are added; the exact cardinality $2r+1$ shows why one
plateau cannot provide the required infinite subsequence.\\
Spectral projector identity &
rigorous reduction, not r429 triangle &
It retains signed eigenvector alignment.  Without a source estimate
for the projector mass it remains a reduction, not a proof.\\
Potapov, de Branges, Abel, PNT-density replacement &
not repeated &
No product index, phase dominance, left--right split, or replacement
density enters the argument.\\
\bottomrule
\end{tabular}
\end{center}

\section*{Claim boundary}

This note proves finite linear-algebra statements and reports a
finite full-comb census.  It does not prove the race lemma for an
infinite subsequence.  Nothing here is evidence for or against the
Riemann Hypothesis in either direction.  No RH claim.  No anti-RH
claim.

\end{document}
